%====================================================%
% SEZIONE 1: Definizione
%====================================================%

\begin{question}[ID=1.1, subtitle={Classificazione di equazioni}, difficulty=facile, tags={equazioni}]
Stabilisci se le seguenti equazioni sono determinate, indeterminate o impossibili.
\begin{enumerate}[leftmargin=3cm]
    \item $\makebox[3cm][l]{5x = 0}\rule{5cm}{0.2pt}$
    \item $\makebox[3cm][l]{4x = 4x + 1}\rule{5cm}{0.2pt}$
    \item $\makebox[3cm][l]{2(x+1) - 2 = 2x}\rule{5cm}{0.2pt}$
\end{enumerate}
\begin{solution}
\begin{enumerate}
    \item determinata ($x=0$)
    \item impossibile
    \item indeterminata (infinite soluzioni)
\end{enumerate}
\end{solution}
\begin{explanation}
Si porta l'equazione nella forma $ax=b$ e si guarda il coefficiente $a$:
se $a\neq 0$ l'equazione è \textbf{determinata} (una sola soluzione); se $a=0$ e $b=0$ è \textbf{indeterminata} (ogni numero è soluzione); se $a=0$ e $b\neq 0$ è \textbf{impossibile}.
\begin{enumerate}
    \item \[\begin{aligned}
    5x &= 0\\
    x &= \dfrac{0}{5}=0 && \textcolor{gray}{\text{divido per $5$}}
    \end{aligned}\]
    Il coefficiente di $x$ è $5\neq0$: determinata.
    \item \[\begin{aligned}
    4x &= 4x+1\\
    4x-4x &= 1 && \textcolor{gray}{\text{porto le $x$ al primo membro}}\\
    0\cdot x &= 1
    \end{aligned}\]
    Nessun numero moltiplicato per $0$ dà $1$: impossibile.
    \item \[\begin{aligned}
    2(x+1)-2 &= 2x\\
    2x+2-2 &= 2x && \textcolor{gray}{\text{svolgo la parentesi}}\\
    2x-2x &= 0 && \textcolor{gray}{\text{porto le $x$ al primo membro}}\\
    0\cdot x &= 0
    \end{aligned}\]
    L'uguaglianza è vera per ogni $x$: indeterminata.
\end{enumerate}
\end{explanation}
\end{question}

%----------------------------------------------------%

\begin{question}[ID=1.2, subtitle={Classificazione di equazioni}, difficulty=facile, tags={equazioni}]
Stabilisci se le seguenti equazioni sono determinate, indeterminate o impossibili.
\begin{enumerate}[leftmargin=3cm]
    \item $\makebox[3cm][l]{7x - 14 = 0}\rule{5cm}{0.2pt}$
    \item $\makebox[3cm][l]{3x + 2 = 3x + 2}\rule{5cm}{0.2pt}$
    \item $\makebox[3cm][l]{5x - 1 = 5x + 3}\rule{5cm}{0.2pt}$
\end{enumerate}
\begin{solution}
\begin{enumerate}
    \item determinata ($x=2$)
    \item indeterminata (infinite soluzioni)
    \item impossibile
\end{enumerate}
\end{solution}
\begin{explanation}
\begin{enumerate}
    \item \[\begin{aligned}
    7x-14 &= 0\\
    7x &= 14 && \textcolor{gray}{\text{sposto il $-14$ cambiando segno}}\\
    x &= 2 && \textcolor{gray}{\text{divido per $7$}}
    \end{aligned}\]
    Una sola soluzione: determinata.
    \item \[\begin{aligned}
    3x+2 &= 3x+2\\
    3x-3x &= 2-2 && \textcolor{gray}{\text{porto le $x$ a sinistra, i numeri a destra}}\\
    0\cdot x &= 0
    \end{aligned}\]
    Vera per ogni $x$: indeterminata.
    \item \[\begin{aligned}
    5x-1 &= 5x+3\\
    5x-5x &= 3+1\\
    0\cdot x &= 4
    \end{aligned}\]
    Mai vera: impossibile.
\end{enumerate}
\end{explanation}
\end{question}

%----------------------------------------------------%

\begin{question}[ID=1.3, subtitle={Classificazione di equazioni}, difficulty=facile, tags={equazioni}]
Stabilisci se le seguenti equazioni sono determinate, indeterminate o impossibili.
\begin{enumerate}[leftmargin=3cm]
    \item $\makebox[3cm][l]{2(x+3) = 2x + 6}\rule{5cm}{0.2pt}$
    \item $\makebox[3cm][l]{4x - 5 = 4x - 8}\rule{5cm}{0.2pt}$
    \item $\makebox[3cm][l]{3(x-2) = 3x - 5}\rule{5cm}{0.2pt}$
\end{enumerate}
\begin{solution}
\begin{enumerate}
    \item indeterminata (infinite soluzioni)
    \item impossibile
    \item impossibile
\end{enumerate}
\end{solution}
\begin{explanation}
\begin{enumerate}
    \item \[\begin{aligned}
    2(x+3) &= 2x+6\\
    2x+6 &= 2x+6 && \textcolor{gray}{\text{svolgo la parentesi}}\\
    0\cdot x &= 0 && \textcolor{gray}{\text{porto tutto al primo membro}}
    \end{aligned}\]
    Vera per ogni $x$: indeterminata.
    \item \[\begin{aligned}
    4x-5 &= 4x-8\\
    4x-4x &= -8+5\\
    0\cdot x &= -3
    \end{aligned}\]
    Mai vera: impossibile.
    \item \[\begin{aligned}
    3(x-2) &= 3x-5\\
    3x-6 &= 3x-5 && \textcolor{gray}{\text{svolgo la parentesi}}\\
    0\cdot x &= 1
    \end{aligned}\]
    Mai vera: impossibile.
\end{enumerate}
\end{explanation}
\end{question}

%----------------------------------------------------%

\begin{question}[ID=1.4, subtitle={Classificazione di equazioni con parentesi}, difficulty=media, tags={equazioni}]
Stabilisci se le seguenti equazioni sono determinate, indeterminate o impossibili.
\begin{enumerate}[leftmargin=3cm]
    \item $\makebox[3cm][l]{3(2x-1) = 6x - 3}\rule{5cm}{0.2pt}$
    \item $\makebox[3cm][l]{5(x+2) - 3x = 2x + 7}\rule{5cm}{0.2pt}$
    \item $\makebox[3cm][l]{4(2x-1) = 8x + 3}\rule{5cm}{0.2pt}$
\end{enumerate}
\begin{solution}
\begin{enumerate}
    \item indeterminata (infinite soluzioni)
    \item impossibile
    \item impossibile
\end{enumerate}
\end{solution}
\begin{explanation}
\begin{enumerate}
    \item \[\begin{aligned}
    3(2x-1) &= 6x-3\\
    6x-3 &= 6x-3 && \textcolor{gray}{\text{svolgo la parentesi}}\\
    0\cdot x &= 0
    \end{aligned}\]
    I due membri sono identici: indeterminata.
    \item \[\begin{aligned}
    5(x+2)-3x &= 2x+7\\
    5x+10-3x &= 2x+7 && \textcolor{gray}{\text{svolgo la parentesi}}\\
    2x+10 &= 2x+7 && \textcolor{gray}{\text{sommo i termini simili}}\\
    0\cdot x &= -3
    \end{aligned}\]
    Mai vera: impossibile.
    \item \[\begin{aligned}
    4(2x-1) &= 8x+3\\
    8x-4 &= 8x+3\\
    0\cdot x &= 7
    \end{aligned}\]
    Mai vera: impossibile.
\end{enumerate}
\end{explanation}
\end{question}

%----------------------------------------------------%

\begin{question}[ID=1.5, subtitle={Classificazione di equazioni con frazioni}, difficulty=difficile, tags={equazioni}]
Stabilisci se le seguenti equazioni sono determinate, indeterminate o impossibili.
\begin{enumerate}[leftmargin=3cm]
    \item $\makebox[3cm][l]{\dfrac{3x-2}{2} = \dfrac{6x-4}{4}}\rule{5cm}{0.2pt}$
    \item $\makebox[3cm][l]{2(3x-1) - (4x+5) = 2x - 7}\rule{5cm}{0.2pt}$
    \item $\makebox[3cm][l]{\dfrac{2x+3}{3} + \dfrac{x-1}{2} = \dfrac{7x+4}{6}}\rule{5cm}{0.2pt}$
\end{enumerate}
\begin{solution}
\begin{enumerate}
    \item indeterminata (infinite soluzioni)
    \item indeterminata (infinite soluzioni)
    \item impossibile
\end{enumerate}
\end{solution}
\begin{explanation}
\begin{enumerate}
    \item \[\begin{aligned}
    \dfrac{3x-2}{2} &= \dfrac{6x-4}{4}\\
    2(3x-2) &= 6x-4 && \textcolor{gray}{\text{moltiplico per il m.c.m. $4$}}\\
    6x-4 &= 6x-4\\
    0\cdot x &= 0
    \end{aligned}\]
    Indeterminata (infatti $\tfrac{6x-4}{4}$ semplificata è proprio $\tfrac{3x-2}{2}$).
    \item \[\begin{aligned}
    2(3x-1)-(4x+5) &= 2x-7\\
    6x-2-4x-5 &= 2x-7 && \textcolor{gray}{\text{svolgo le parentesi}}\\
    2x-7 &= 2x-7\\
    0\cdot x &= 0
    \end{aligned}\]
    Indeterminata.
    \item \[\begin{aligned}
    \dfrac{2x+3}{3}+\dfrac{x-1}{2} &= \dfrac{7x+4}{6}\\
    2(2x+3)+3(x-1) &= 7x+4 && \textcolor{gray}{\text{moltiplico per il m.c.m. $6$}}\\
    4x+6+3x-3 &= 7x+4 && \textcolor{gray}{\text{svolgo le parentesi}}\\
    7x+3 &= 7x+4\\
    0\cdot x &= 1
    \end{aligned}\]
    Mai vera: impossibile.
\end{enumerate}
\end{explanation}
\end{question}


%====================================================%
% SEZIONE 2: Equazioni lineari
%====================================================%

\begin{question}[ID=2.1, subtitle={Equazioni di primo grado}, difficulty=facile, tags={equazioni lineari}]
Risolvi le seguenti equazioni di primo grado:
\begin{multicols}{2}
\begin{enumerate}[leftmargin=3cm]
    \item $x + 5 = 7$
    \item $3x - 7 = 2x + 4$
    \item $2(x - 3) + 4 = 3x - 2$
    \item $\dfrac{1}{2}x + 1 = \dfrac{2}{3}(x + 1)$
\end{enumerate}
\end{multicols}
\begin{solution}
\begin{enumerate}
    \item $x=2$
    \item $x=11$
    \item $x=0$
    \item $x=2$
\end{enumerate}
\end{solution}
\begin{explanation}
Si isola la $x$: si svolgono le parentesi, si portano i termini con la $x$ al primo membro e i numeri al secondo (cambiando segno a ciò che si sposta), poi si divide per il coefficiente della $x$.
\begin{enumerate}
    \item \[\begin{aligned}
    x+5 &= 7\\
    x &= 7-5 && \textcolor{gray}{\text{sposto il $5$ cambiando segno}}\\
    x &= 2
    \end{aligned}\]
    \item \[\begin{aligned}
    3x-7 &= 2x+4\\
    3x-2x &= 4+7 && \textcolor{gray}{\text{$x$ a sinistra, numeri a destra}}\\
    x &= 11
    \end{aligned}\]
    \item \[\begin{aligned}
    2(x-3)+4 &= 3x-2\\
    2x-6+4 &= 3x-2 && \textcolor{gray}{\text{svolgo la parentesi}}\\
    2x-3x &= -2+2 && \textcolor{gray}{\text{$x$ a sinistra, numeri a destra}}\\
    -x &= 0\\
    x &= 0
    \end{aligned}\]
    \item \[\begin{aligned}
    \dfrac12x+1 &= \dfrac23(x+1)\\
    3x+6 &= 4(x+1) && \textcolor{gray}{\text{moltiplico per il m.c.m. $6$}}\\
    3x+6 &= 4x+4 && \textcolor{gray}{\text{svolgo la parentesi}}\\
    3x-4x &= 4-6\\
    -x &= -2\\
    x &= 2 && \textcolor{gray}{\text{moltiplico per $-1$}}
    \end{aligned}\]
\end{enumerate}
\end{explanation}
\end{question}

%----------------------------------------------------%

\begin{question}[ID=2.2, subtitle={Equazioni di primo grado}, difficulty=facile, tags={equazioni lineari}]
Risolvi le seguenti equazioni di primo grado:
\begin{multicols}{2}
\begin{enumerate}[leftmargin=3cm]
    \item $2x + 3 = 11$
    \item $5x - 4 = 2x + 8$
    \item $3(x + 2) - 5 = 2x + 4$
    \item $\dfrac{3}{4}x - 2 = \dfrac{1}{2}x + 1$
\end{enumerate}
\end{multicols}
\begin{solution}
\begin{enumerate}
    \item $x=4$
    \item $x=4$
    \item $x=3$
    \item $x=12$
\end{enumerate}
\end{solution}
\begin{explanation}
\begin{enumerate}
    \item \[\begin{aligned}
    2x+3 &= 11\\
    2x &= 8 && \textcolor{gray}{\text{sposto il $3$}}\\
    x &= 4 && \textcolor{gray}{\text{divido per $2$}}
    \end{aligned}\]
    \item \[\begin{aligned}
    5x-4 &= 2x+8\\
    5x-2x &= 8+4 && \textcolor{gray}{\text{$x$ a sinistra, numeri a destra}}\\
    3x &= 12\\
    x &= 4 && \textcolor{gray}{\text{divido per $3$}}
    \end{aligned}\]
    \item \[\begin{aligned}
    3(x+2)-5 &= 2x+4\\
    3x+6-5 &= 2x+4 && \textcolor{gray}{\text{svolgo la parentesi}}\\
    3x-2x &= 4-1\\
    x &= 3
    \end{aligned}\]
    \item \[\begin{aligned}
    \dfrac34x-2 &= \dfrac12x+1\\
    3x-8 &= 2x+4 && \textcolor{gray}{\text{moltiplico per il m.c.m. $4$}}\\
    3x-2x &= 4+8\\
    x &= 12
    \end{aligned}\]
\end{enumerate}
\end{explanation}
\end{question}

%----------------------------------------------------%

\begin{question}[ID=2.3, subtitle={Equazioni di primo grado}, difficulty=facile, tags={equazioni lineari}]
Risolvi le seguenti equazioni di primo grado:
\begin{multicols}{2}
\begin{enumerate}[leftmargin=3cm]
    \item $4x + 1 = 2x + 9$
    \item $7 - 3x = x - 9$
    \item $2(3x - 1) + 4 = 5x + 7$
    \item $\dfrac{2x+1}{3} = \dfrac{x-2}{2}$
\end{enumerate}
\end{multicols}
\begin{solution}
\begin{enumerate}
    \item $x=4$
    \item $x=4$
    \item $x=5$
    \item $x=-8$
\end{enumerate}
\end{solution}
\begin{explanation}
\begin{enumerate}
    \item \[\begin{aligned}
    4x+1 &= 2x+9\\
    4x-2x &= 9-1 && \textcolor{gray}{\text{$x$ a sinistra, numeri a destra}}\\
    2x &= 8\\
    x &= 4 && \textcolor{gray}{\text{divido per $2$}}
    \end{aligned}\]
    \item \[\begin{aligned}
    7-3x &= x-9\\
    -3x-x &= -9-7\\
    -4x &= -16\\
    x &= 4 && \textcolor{gray}{\text{divido per $-4$}}
    \end{aligned}\]
    \item \[\begin{aligned}
    2(3x-1)+4 &= 5x+7\\
    6x-2+4 &= 5x+7 && \textcolor{gray}{\text{svolgo la parentesi}}\\
    6x-5x &= 7-2\\
    x &= 5
    \end{aligned}\]
    \item \[\begin{aligned}
    \dfrac{2x+1}{3} &= \dfrac{x-2}{2}\\
    2(2x+1) &= 3(x-2) && \textcolor{gray}{\text{moltiplico per il m.c.m. $6$}}\\
    4x+2 &= 3x-6\\
    4x-3x &= -6-2\\
    x &= -8
    \end{aligned}\]
\end{enumerate}
\end{explanation}
\end{question}

%----------------------------------------------------%

\begin{question}[ID=2.4, subtitle={Equazioni di primo grado}, difficulty=media, tags={equazioni lineari}]
Risolvi le seguenti equazioni di primo grado:
\begin{multicols}{2}
\begin{enumerate}[leftmargin=3cm]
    \item $4(2x-3) - 5 = 3x + 6$
    \item $7 - 2(3x+1) = 4x - 9$
    \item $\dfrac{3x-2}{4} + \dfrac{x+1}{2} = 5$
    \item $5(x-2) - 2(2x+1) = 3 - x$
\end{enumerate}
\end{multicols}
\begin{solution}
\begin{enumerate}
    \item $x=\tfrac{23}{5}$
    \item $x=\tfrac{7}{5}$
    \item $x=4$
    \item $x=\tfrac{15}{2}$
\end{enumerate}
\end{solution}
\begin{explanation}
\begin{enumerate}
    \item \[\begin{aligned}
    4(2x-3)-5 &= 3x+6\\
    8x-12-5 &= 3x+6 && \textcolor{gray}{\text{svolgo la parentesi}}\\
    8x-3x &= 6+17 && \textcolor{gray}{\text{$x$ a sinistra, numeri a destra}}\\
    5x &= 23\\
    x &= \dfrac{23}{5} && \textcolor{gray}{\text{divido per $5$}}
    \end{aligned}\]
    \item \[\begin{aligned}
    7-2(3x+1) &= 4x-9\\
    7-6x-2 &= 4x-9 && \textcolor{gray}{\text{svolgo la parentesi (il meno cambia i segni)}}\\
    -6x-4x &= -9-5\\
    -10x &= -14\\
    x &= \dfrac{14}{10}=\dfrac75 && \textcolor{gray}{\text{divido per $-10$ e semplifico}}
    \end{aligned}\]
    \item \[\begin{aligned}
    \dfrac{3x-2}{4}+\dfrac{x+1}{2} &= 5\\
    3x-2+2(x+1) &= 20 && \textcolor{gray}{\text{moltiplico per il m.c.m. $4$}}\\
    3x-2+2x+2 &= 20\\
    5x &= 20\\
    x &= 4
    \end{aligned}\]
    \item \[\begin{aligned}
    5(x-2)-2(2x+1) &= 3-x\\
    5x-10-4x-2 &= 3-x && \textcolor{gray}{\text{svolgo le parentesi}}\\
    x-12 &= 3-x\\
    x+x &= 3+12\\
    2x &= 15\\
    x &= \dfrac{15}{2}
    \end{aligned}\]
\end{enumerate}
\end{explanation}
\end{question}

%----------------------------------------------------%

\begin{question}[ID=2.5, subtitle={Equazioni di primo grado con frazioni}, difficulty=difficile, tags={equazioni lineari}]
Risolvi le seguenti equazioni di primo grado:
\begin{multicols}{2}
\begin{enumerate}[leftmargin=3cm]
    \item $\dfrac{2x-1}{3} - \dfrac{x+2}{4} = \dfrac{5x-7}{6}$
    \item $2\left(\dfrac{x-3}{4}\right) - \dfrac{3x+1}{5} = \dfrac{x-7}{10}$
    \item $\dfrac{3}{2}(x-4) - \dfrac{2}{3}(x+3) = \dfrac{x-5}{6}$
    \item $\dfrac{5x-2}{6} - \dfrac{2x+1}{3} + \dfrac{x-4}{2} = 7$
\end{enumerate}
\end{multicols}
\begin{solution}
\begin{enumerate}
    \item $x=\tfrac45$
    \item $x=-5$
    \item $x=\tfrac{43}{4}$
    \item $x=\tfrac{29}{2}$
\end{enumerate}
\end{solution}
\begin{explanation}
Si moltiplicano entrambi i membri per il m.c.m. dei denominatori, così l'equazione diventa senza frazioni.
\begin{enumerate}
    \item \[\begin{aligned}
    \dfrac{2x-1}{3}-\dfrac{x+2}{4} &= \dfrac{5x-7}{6}\\
    4(2x-1)-3(x+2) &= 2(5x-7) && \textcolor{gray}{\text{moltiplico per il m.c.m. $12$}}\\
    8x-4-3x-6 &= 10x-14 && \textcolor{gray}{\text{svolgo le parentesi}}\\
    5x-10 &= 10x-14\\
    5x-10x &= -14+10\\
    -5x &= -4\\
    x &= \dfrac45
    \end{aligned}\]
    \item Prima semplifico: $2\cdot\dfrac{x-3}{4}=\dfrac{x-3}{2}$.
    \[\begin{aligned}
    \dfrac{x-3}{2}-\dfrac{3x+1}{5} &= \dfrac{x-7}{10}\\
    5(x-3)-2(3x+1) &= x-7 && \textcolor{gray}{\text{moltiplico per il m.c.m. $10$}}\\
    5x-15-6x-2 &= x-7\\
    -x-17 &= x-7\\
    -2x &= 10\\
    x &= -5
    \end{aligned}\]
    \item \[\begin{aligned}
    \dfrac32(x-4)-\dfrac23(x+3) &= \dfrac{x-5}{6}\\
    9(x-4)-4(x+3) &= x-5 && \textcolor{gray}{\text{moltiplico per il m.c.m. $6$}}\\
    9x-36-4x-12 &= x-5\\
    5x-48 &= x-5\\
    4x &= 43\\
    x &= \dfrac{43}{4}
    \end{aligned}\]
    \item \[\begin{aligned}
    \dfrac{5x-2}{6}-\dfrac{2x+1}{3}+\dfrac{x-4}{2} &= 7\\
    (5x-2)-2(2x+1)+3(x-4) &= 42 && \textcolor{gray}{\text{moltiplico per il m.c.m. $6$}}\\
    5x-2-4x-2+3x-12 &= 42\\
    4x-16 &= 42\\
    4x &= 58\\
    x &= \dfrac{29}{2} && \textcolor{gray}{\text{divido per $4$ e semplifico}}
    \end{aligned}\]
\end{enumerate}
\end{explanation}
\end{question}


%====================================================%
% SEZIONE 3: Equazioni frazionarie
%====================================================%

\begin{question}[ID=3.1, subtitle={Equazioni frazionarie}, difficulty=facile, tags={equazioni frazionarie}]
Risolvi le seguenti equazioni frazionarie:
\begin{multicols}{2}
\begin{enumerate}[leftmargin=3cm]
    \item $\dfrac{x+1}{x-2} = 0$
    \item $\dfrac{(x-3) \cdot (x+3)}{2x+6} = 0$
    \item $\dfrac{2x-1}{3x-1} = 2$
    \item $\dfrac{3}{x-3} = \dfrac{2}{2x+4}$
    \item $\dfrac{1+3x}{4x+4} - \dfrac{5-x}{x+1} = 2$
\end{enumerate}
\end{multicols}
\begin{solution}
\begin{enumerate}
    \item $x=-1$
    \item $x=3$
    \item $x=\tfrac14$
    \item $x=-\tfrac92$
    \item $x=-27$
\end{enumerate}
\end{solution}
\begin{explanation}
In un'equazione fratta si fissano prima le \textbf{condizioni di esistenza} (C.E.): i denominatori non possono valere zero. Poi si eliminano i denominatori moltiplicando per il loro m.c.m. e, alla fine, si controlla che la soluzione rispetti le C.E.
\begin{enumerate}
    \item C.E.: $x-2\neq0$, cioè $x\neq2$. \\
    Una frazione vale zero solo quando il numeratore è zero:
    \[ x+1=0 \;\Rightarrow\; x=-1 \]
    $-1\neq2$: soluzione accettabile.
    \item C.E.: $2x+6\neq0$, cioè $x\neq-3$. \\
    Il numeratore si annulla se uno dei due fattori è zero:
    \[ (x-3)(x+3)=0 \;\Rightarrow\; x=3 \;\text{ oppure }\; x=-3 \]
    $x=-3$ non rispetta le C.E. e va scartata; resta $x=3$.
    \item C.E.: $3x-1\neq0$, cioè $x\neq\tfrac13$.
    \[\begin{aligned}
    \dfrac{2x-1}{3x-1} &= 2\\
    2x-1 &= 2(3x-1) && \textcolor{gray}{\text{moltiplico per $3x-1$}}\\
    2x-1 &= 6x-2\\
    -4x &= -1\\
    x &= \dfrac14
    \end{aligned}\]
    $\tfrac14\neq\tfrac13$: accettabile.
    \item C.E.: $x\neq3$ e $x\neq-2$.
    \[\begin{aligned}
    \dfrac{3}{x-3} &= \dfrac{2}{2x+4}\\
    3(2x+4) &= 2(x-3) && \textcolor{gray}{\text{moltiplico in croce}}\\
    6x+12 &= 2x-6\\
    4x &= -18\\
    x &= -\dfrac92
    \end{aligned}\]
    Rispetta le C.E.: accettabile.
    \item C.E.: $4x+4=4(x+1)\neq0$, cioè $x\neq-1$.
    \[\begin{aligned}
    \dfrac{1+3x}{4(x+1)}-\dfrac{5-x}{x+1} &= 2\\
    1+3x-4(5-x) &= 8(x+1) && \textcolor{gray}{\text{moltiplico per il m.c.m. $4(x+1)$}}\\
    1+3x-20+4x &= 8x+8\\
    7x-19 &= 8x+8\\
    -x &= 27\\
    x &= -27
    \end{aligned}\]
    Rispetta le C.E.: accettabile.
\end{enumerate}
\end{explanation}
\end{question}

%----------------------------------------------------%

\begin{question}[ID=3.2, subtitle={Equazioni frazionarie}, difficulty=facile, tags={equazioni frazionarie}]
Risolvi le seguenti equazioni frazionarie:
\begin{multicols}{2}
\begin{enumerate}[leftmargin=3cm]
    \item $\dfrac{x-4}{x+1} = 0$
    \item $\dfrac{(x-2)(x+2)}{x-2} = 0$
    \item $\dfrac{x+3}{2x-1} = 1$
    \item $\dfrac{4}{x+2} = \dfrac{2}{x-1}$
    \item $\dfrac{2x+1}{x-2} - \dfrac{x+4}{x-2} = 3$
\end{enumerate}
\end{multicols}
\begin{solution}
\begin{enumerate}
    \item $x=4$
    \item $x=-2$
    \item $x=4$
    \item $x=4$
    \item $x=\tfrac32$
\end{enumerate}
\end{solution}
\begin{explanation}
\begin{enumerate}
    \item C.E.: $x\neq-1$. \ Il numeratore deve essere zero: $x-4=0 \Rightarrow x=4$, accettabile.
    \item C.E.: $x\neq2$. \ Il numeratore deve essere zero:
    \[ (x-2)(x+2)=0 \;\Rightarrow\; x=2 \;\text{ oppure }\; x=-2 \]
    $x=2$ è escluso dalle C.E.; resta $x=-2$.
    \item C.E.: $x\neq\tfrac12$.
    \[\begin{aligned}
    x+3 &= 2x-1 && \textcolor{gray}{\text{moltiplico per $2x-1$}}\\
    -x &= -4\\
    x &= 4
    \end{aligned}\]
    Accettabile.
    \item C.E.: $x\neq-2$ e $x\neq1$.
    \[\begin{aligned}
    4(x-1) &= 2(x+2) && \textcolor{gray}{\text{moltiplico in croce}}\\
    4x-4 &= 2x+4\\
    2x &= 8\\
    x &= 4
    \end{aligned}\]
    Accettabile.
    \item C.E.: $x\neq2$. \ Le due frazioni hanno lo stesso denominatore, quindi sottraggo i numeratori:
    \[\begin{aligned}
    \dfrac{(2x+1)-(x+4)}{x-2} &= 3\\
    \dfrac{x-3}{x-2} &= 3\\
    x-3 &= 3(x-2) && \textcolor{gray}{\text{moltiplico per $x-2$}}\\
    x-3 &= 3x-6\\
    -2x &= -3\\
    x &= \dfrac32
    \end{aligned}\]
    Accettabile.
\end{enumerate}
\end{explanation}
\end{question}

%----------------------------------------------------%

\begin{question}[ID=3.3, subtitle={Equazioni frazionarie}, difficulty=facile, tags={equazioni frazionarie}]
Risolvi le seguenti equazioni frazionarie:
\begin{multicols}{2}
\begin{enumerate}[leftmargin=3cm]
    \item $\dfrac{2x-6}{x+2} = 0$
    \item $\dfrac{(x+1)(x-4)}{x+1} = 0$
    \item $\dfrac{3x+2}{x-1} = 2$
    \item $\dfrac{5}{x+1} = \dfrac{3}{x-1}$
    \item $\dfrac{x+2}{3x-3} + \dfrac{2x-1}{3x-3} = 2$
\end{enumerate}
\end{multicols}
\begin{solution}
\begin{enumerate}
    \item $x=3$
    \item $x=4$
    \item $x=-4$
    \item $x=4$
    \item $x=\tfrac73$
\end{enumerate}
\end{solution}
\begin{explanation}
\begin{enumerate}
    \item C.E.: $x\neq-2$. \ Numeratore nullo: $2x-6=0 \Rightarrow x=3$, accettabile.
    \item C.E.: $x\neq-1$. \ Numeratore nullo: $(x+1)(x-4)=0 \Rightarrow x=-1$ oppure $x=4$. \\
    $x=-1$ è escluso dalle C.E.; resta $x=4$.
    \item C.E.: $x\neq1$.
    \[\begin{aligned}
    3x+2 &= 2(x-1) && \textcolor{gray}{\text{moltiplico per $x-1$}}\\
    3x+2 &= 2x-2\\
    x &= -4
    \end{aligned}\]
    Accettabile.
    \item C.E.: $x\neq\pm1$.
    \[\begin{aligned}
    5(x-1) &= 3(x+1) && \textcolor{gray}{\text{moltiplico in croce}}\\
    5x-5 &= 3x+3\\
    2x &= 8\\
    x &= 4
    \end{aligned}\]
    Accettabile.
    \item C.E.: $3x-3\neq0$, cioè $x\neq1$. \ Stesso denominatore: sommo i numeratori.
    \[\begin{aligned}
    \dfrac{x+2+2x-1}{3x-3} &= 2\\
    3x+1 &= 2(3x-3) && \textcolor{gray}{\text{moltiplico per $3x-3$}}\\
    3x+1 &= 6x-6\\
    -3x &= -7\\
    x &= \dfrac73
    \end{aligned}\]
    Accettabile.
\end{enumerate}
\end{explanation}
\end{question}

%----------------------------------------------------%

\begin{question}[ID=3.4, subtitle={Equazioni frazionarie}, difficulty=media, tags={equazioni frazionarie}]
Risolvi le seguenti equazioni frazionarie:
\begin{multicols}{2}
\begin{enumerate}[leftmargin=3cm]
    \item $\dfrac{x-5}{x+2} = 2$
    \item $\dfrac{2x+1}{x-3} = \dfrac{3}{x-3}$
    \item $\dfrac{x+4}{x-1} - 1 = 0$
    \item $\dfrac{3}{x+1} = \dfrac{x}{x+1}$
    \item $\dfrac{2x-3}{x+2} + \dfrac{x+1}{x+2} = 4$
\end{enumerate}
\end{multicols}
\begin{solution}
\begin{enumerate}
    \item $x=-9$
    \item $x=1$
    \item impossibile
    \item $x=3$
    \item $x=-10$
\end{enumerate}
\end{solution}
\begin{explanation}
\begin{enumerate}
    \item C.E.: $x\neq-2$.
    \[\begin{aligned}
    x-5 &= 2(x+2) && \textcolor{gray}{\text{moltiplico per $x+2$}}\\
    x-5 &= 2x+4\\
    -x &= 9\\
    x &= -9
    \end{aligned}\]
    Accettabile.
    \item C.E.: $x\neq3$. \ Le due frazioni hanno lo stesso denominatore, quindi devono avere lo stesso numeratore:
    \[ 2x+1=3 \;\Rightarrow\; x=1 \]
    Accettabile.
    \item C.E.: $x\neq1$.
    \[\begin{aligned}
    \dfrac{x+4}{x-1}-1 &= 0\\
    \dfrac{x+4-(x-1)}{x-1} &= 0 && \textcolor{gray}{\text{denominatore comune}}\\
    \dfrac{5}{x-1} &= 0
    \end{aligned}\]
    Il numeratore vale sempre $5$ e non si annulla mai: l'equazione è \textbf{impossibile}.
    \item C.E.: $x\neq-1$. \ Stesso denominatore, quindi stesso numeratore: $x=3$, accettabile.
    \item C.E.: $x\neq-2$. \ Stesso denominatore: sommo i numeratori.
    \[\begin{aligned}
    \dfrac{2x-3+x+1}{x+2} &= 4\\
    3x-2 &= 4(x+2) && \textcolor{gray}{\text{moltiplico per $x+2$}}\\
    3x-2 &= 4x+8\\
    -x &= 10\\
    x &= -10
    \end{aligned}\]
    Accettabile.
\end{enumerate}
\end{explanation}
\end{question}

%----------------------------------------------------%

\begin{question}[ID=3.5, subtitle={Equazioni frazionarie con più denominatori}, difficulty=difficile, tags={equazioni frazionarie}]
Risolvi le seguenti equazioni frazionarie:
\begin{multicols}{2}
\begin{enumerate}[leftmargin=3cm]
    \item $\dfrac{2x+1}{x-2} - \dfrac{x-3}{x+1} = 1$
    \item $\dfrac{3}{x-1} + \dfrac{2}{x+2} = \dfrac{5x+1}{(x-1)(x+2)}$
    \item $\dfrac{x+2}{x-3} = \dfrac{2x-1}{x+1}$
    \item $\dfrac{4}{x-2} - \dfrac{1}{x+2} = \dfrac{3x+10}{x^2-4}$
    \item $\dfrac{x-1}{x+3} + \dfrac{x+2}{x-2} = 2$
\end{enumerate}
\end{multicols}
\begin{solution}
\begin{enumerate}
    \item $x=\tfrac13$
    \item impossibile
    \item $x=5\pm2\sqrt6$
    \item indeterminata: ogni $x\neq\pm2$
    \item impossibile
\end{enumerate}
\end{solution}
\begin{explanation}
\begin{enumerate}
    \item C.E.: $x\neq2$ e $x\neq-1$. Moltiplico per il m.c.m. $(x-2)(x+1)$:
    \[\begin{aligned}
    (2x+1)(x+1)-(x-3)(x-2) &= (x-2)(x+1)\\
    2x^2+3x+1-(x^2-5x+6) &= x^2-x-2 && \textcolor{gray}{\text{sviluppo i prodotti}}\\
    x^2+8x-5 &= x^2-x-2\\
    9x &= 3 && \textcolor{gray}{\text{le $x^2$ si eliminano}}\\
    x &= \dfrac13
    \end{aligned}\]
    Accettabile.
    \item C.E.: $x\neq1$ e $x\neq-2$. Moltiplico per il m.c.m. $(x-1)(x+2)$:
    \[\begin{aligned}
    3(x+2)+2(x-1) &= 5x+1\\
    3x+6+2x-2 &= 5x+1\\
    5x+4 &= 5x+1\\
    0\cdot x &= -3
    \end{aligned}\]
    Mai vera: impossibile.
    \item C.E.: $x\neq3$ e $x\neq-1$.
    \[\begin{aligned}
    (x+2)(x+1) &= (2x-1)(x-3) && \textcolor{gray}{\text{moltiplico in croce}}\\
    x^2+3x+2 &= 2x^2-7x+3 && \textcolor{gray}{\text{sviluppo i prodotti}}\\
    0 &= x^2-10x+1 && \textcolor{gray}{\text{porto tutto a destra}}
    \end{aligned}\]
    Equazione di secondo grado: $x=\dfrac{10\pm\sqrt{100-4}}{2}=\dfrac{10\pm\sqrt{96}}{2}=\dfrac{10\pm4\sqrt6}{2}=5\pm2\sqrt6$. \\
    Nessuna delle due vale $3$ o $-1$: entrambe accettabili.
    \item Poiché $x^2-4=(x-2)(x+2)$, le C.E. sono $x\neq\pm2$ e il m.c.m. è $x^2-4$.
    \[\begin{aligned}
    4(x+2)-(x-2) &= 3x+10\\
    4x+8-x+2 &= 3x+10\\
    3x+10 &= 3x+10
    \end{aligned}\]
    L'uguaglianza è sempre vera: l'equazione è \textbf{indeterminata}, cioè ogni $x$ che rispetta le C.E. è soluzione.
    \item C.E.: $x\neq-3$ e $x\neq2$. Moltiplico per il m.c.m. $(x+3)(x-2)$:
    \[\begin{aligned}
    (x-1)(x-2)+(x+2)(x+3) &= 2(x+3)(x-2)\\
    x^2-3x+2+x^2+5x+6 &= 2(x^2+x-6) && \textcolor{gray}{\text{sviluppo i prodotti}}\\
    2x^2+2x+8 &= 2x^2+2x-12\\
    0\cdot x &= -20
    \end{aligned}\]
    Mai vera: impossibile.
\end{enumerate}
\end{explanation}
\end{question}

%====================================================%
% SEZIONE 4: Equazioni di grado superiore al secondo
%====================================================%

\begin{question}[ID=4.1, subtitle={Equazioni di grado superiore al secondo}, difficulty=facile, tags={equazioni di grado superiore}]
Risolvi le seguenti equazioni:
\begin{multicols}{2}
\begin{enumerate}[leftmargin=3cm]
    \item $x^3 - 8 = 0$
    \item $x^3 + 27 = 0$
    \item $x^4 - 16 = 0$
    \item $x^4 - 81 = 0$
    \item $x^3 - 4x = 0$
    \item $x^3 - 9x = 0$
\end{enumerate}
\end{multicols}
\begin{solution}
\begin{enumerate}
    \item $x=2$
    \item $x=-3$
    \item $x=\pm2$
    \item $x=\pm3$
    \item $x=-2,\ 0,\ 2$
    \item $x=-3,\ 0,\ 3$
\end{enumerate}
\end{solution}
\begin{explanation}
\begin{enumerate}
    \item \[\begin{aligned}
    x^3 &= 8\\
    x &= \sqrt[3]{8}=2 && \textcolor{gray}{\text{la radice cubica ha una sola soluzione reale}}
    \end{aligned}\]
    \item \[\begin{aligned}
    x^3 &= -27\\
    x &= \sqrt[3]{-27}=-3 && \textcolor{gray}{\text{la radice cubica di un negativo è negativa}}
    \end{aligned}\]
    \item Scompongo come differenza di quadrati:
    \[\begin{aligned}
    x^4-16 &= 0\\
    (x^2-4)(x^2+4) &= 0\\
    (x-2)(x+2)(x^2+4) &= 0
    \end{aligned}\]
    $x^2+4$ è sempre positivo, quindi restano $x=\pm2$.
    \item \[\begin{aligned}
    (x^2-9)(x^2+9) &= 0 && \textcolor{gray}{\text{differenza di quadrati}}\\
    (x-3)(x+3)(x^2+9) &= 0
    \end{aligned}\]
    $x^2+9>0$, quindi $x=\pm3$.
    \item \[\begin{aligned}
    x^3-4x &= 0\\
    x(x^2-4) &= 0 && \textcolor{gray}{\text{raccolgo $x$}}\\
    x(x-2)(x+2) &= 0 && \textcolor{gray}{\text{differenza di quadrati}}
    \end{aligned}\]
    Annullo ogni fattore: $x=0$, $x=2$, $x=-2$.
    \item \[\begin{aligned}
    x(x^2-9) &= 0 && \textcolor{gray}{\text{raccolgo $x$}}\\
    x(x-3)(x+3) &= 0
    \end{aligned}\]
    $x=-3,\ 0,\ 3$.
\end{enumerate}
\end{explanation}
\end{question}

%----------------------------------------------------%

\begin{question}[ID=4.2, subtitle={Equazioni di grado superiore al secondo}, difficulty=facile, tags={equazioni di grado superiore}]
Risolvi le seguenti equazioni:
\begin{multicols}{2}
\begin{enumerate}[leftmargin=3cm]
    \item $x^3 - 3x^2 = 0$
    \item $x^3 + 2x^2 = 0$
    \item $x^4 - 5x^2 = 0$
    \item $x^4 - 13x^2 + 36 = 0$
    \item $x^3 - 6x^2 + 9x = 0$
    \item $x^4 - 10x^2 + 9 = 0$
\end{enumerate}
\end{multicols}
\begin{solution}
\begin{enumerate}
    \item $x=0,\ 3$
    \item $x=0,\ -2$
    \item $x=0,\ \pm\sqrt5$
    \item $x=\pm2,\ \pm3$
    \item $x=0,\ 3$
    \item $x=\pm1,\ \pm3$
\end{enumerate}
\end{solution}
\begin{explanation}
\begin{enumerate}
    \item Raccolgo $x^2$: \ $x^2(x-3)=0 \;\Rightarrow\; x=0$ (doppia) oppure $x=3$.
    \item Raccolgo $x^2$: \ $x^2(x+2)=0 \;\Rightarrow\; x=0$ oppure $x=-2$.
    \item Raccolgo $x^2$: \ $x^2(x^2-5)=0 \;\Rightarrow\; x=0$ oppure $x^2=5$, cioè $x=\pm\sqrt5$.
    \item È un'equazione \textbf{biquadratica}: pongo $t=x^2$.
    \[\begin{aligned}
    t^2-13t+36 &= 0\\
    t &= \dfrac{13\pm\sqrt{169-144}}{2}=\dfrac{13\pm5}{2} && \textcolor{gray}{\text{formula risolutiva}}\\
    t &= 4 \;\text{ oppure }\; t=9
    \end{aligned}\]
    Torno alla $x$: $x^2=4 \Rightarrow x=\pm2$; \ $x^2=9 \Rightarrow x=\pm3$.
    \item \[\begin{aligned}
    x(x^2-6x+9) &= 0 && \textcolor{gray}{\text{raccolgo $x$}}\\
    x(x-3)^2 &= 0 && \textcolor{gray}{\text{quadrato di binomio}}
    \end{aligned}\]
    $x=0$ oppure $x=3$ (doppia).
    \item Biquadratica, pongo $t=x^2$: \ $t^2-10t+9=0 \Rightarrow (t-1)(t-9)=0 \Rightarrow t=1,\ t=9$. \\
    Quindi $x^2=1 \Rightarrow x=\pm1$; \ $x^2=9 \Rightarrow x=\pm3$.
\end{enumerate}
\end{explanation}
\end{question}

%----------------------------------------------------%

\begin{question}[ID=4.3, subtitle={Equazioni di grado superiore al secondo mediante scomposizione}, difficulty=media, tags={equazioni di grado superiore}]
Risolvi le seguenti equazioni:
\begin{multicols}{2}
\begin{enumerate}[leftmargin=3cm]
    \item $x^3 - 5x^2 + 6x = 0$
    \item $x^3 + x^2 - 6x = 0$
    \item $x^4 - 2x^3 - 8x^2 = 0$
    \item $x^4 - 5x^2 + 4 = 0$
    \item $x^3 - 4x^2 - x + 4 = 0$
    \item $x^3 + 2x^2 - 9x - 18 = 0$
\end{enumerate}
\end{multicols}
\begin{solution}
\begin{enumerate}
    \item $x=0,\ 2,\ 3$
    \item $x=-3,\ 0,\ 2$
    \item $x=-2,\ 0,\ 4$
    \item $x=\pm1,\ \pm2$
    \item $x=-1,\ 1,\ 4$
    \item $x=-3,\ -2,\ 3$
\end{enumerate}
\end{solution}
\begin{explanation}
Si scompone il polinomio in fattori e si applica la legge di annullamento del prodotto.
\begin{enumerate}
    \item \[\begin{aligned}
    x(x^2-5x+6) &= 0 && \textcolor{gray}{\text{raccolgo $x$}}\\
    x(x-2)(x-3) &= 0 && \textcolor{gray}{\text{trinomio: somma $-5$, prodotto $6$}}
    \end{aligned}\]
    $x=0,\ 2,\ 3$.
    \item \[\begin{aligned}
    x(x^2+x-6) &= 0 && \textcolor{gray}{\text{raccolgo $x$}}\\
    x(x+3)(x-2) &= 0 && \textcolor{gray}{\text{somma $1$, prodotto $-6$}}
    \end{aligned}\]
    $x=-3,\ 0,\ 2$.
    \item \[\begin{aligned}
    x^2(x^2-2x-8) &= 0 && \textcolor{gray}{\text{raccolgo $x^2$}}\\
    x^2(x-4)(x+2) &= 0 && \textcolor{gray}{\text{somma $-2$, prodotto $-8$}}
    \end{aligned}\]
    $x=-2,\ 0,\ 4$.
    \item Biquadratica, pongo $t=x^2$: \ $t^2-5t+4=0 \Rightarrow (t-1)(t-4)=0$. \\
    $x^2=1 \Rightarrow x=\pm1$; \ $x^2=4 \Rightarrow x=\pm2$.
    \item Raccoglimento parziale (i primi due termini e gli ultimi due):
    \[\begin{aligned}
    x^2(x-4)-(x-4) &= 0\\
    (x-4)(x^2-1) &= 0 && \textcolor{gray}{\text{raccolgo $(x-4)$}}\\
    (x-4)(x-1)(x+1) &= 0
    \end{aligned}\]
    $x=-1,\ 1,\ 4$.
    \item \[\begin{aligned}
    x^2(x+2)-9(x+2) &= 0 && \textcolor{gray}{\text{raccoglimento parziale}}\\
    (x+2)(x^2-9) &= 0\\
    (x+2)(x-3)(x+3) &= 0
    \end{aligned}\]
    $x=-3,\ -2,\ 3$.
\end{enumerate}
\end{explanation}
\end{question}

%----------------------------------------------------%

\begin{question}[ID=4.4, subtitle={Equazioni di grado superiore al secondo}, difficulty=media, tags={equazioni di grado superiore}]
Risolvi le seguenti equazioni:
\begin{multicols}{2}
\begin{enumerate}[leftmargin=3cm]
    \item $x^4 - 7x^2 + 12 = 0$
    \item $x^4 - 13x^2 + 36 = 0$
    \item $x^4 - 3x^3 - 4x^2 = 0$
    \item $2x^3 - 8x^2 + 6x = 0$
    \item $x^4 - x^3 - 6x^2 = 0$
    \item $x^3 - 7x^2 + 14x - 8 = 0$
\end{enumerate}
\end{multicols}
\begin{solution}
\begin{enumerate}
    \item $x=\pm\sqrt3,\ \pm2$
    \item $x=\pm2,\ \pm3$
    \item $x=-1,\ 0,\ 4$
    \item $x=0,\ 1,\ 3$
    \item $x=-2,\ 0,\ 3$
    \item $x=1,\ 2,\ 4$
\end{enumerate}
\end{solution}
\begin{explanation}
\begin{enumerate}
    \item Biquadratica, pongo $t=x^2$: \ $t^2-7t+12=0 \Rightarrow (t-3)(t-4)=0$. \\
    $x^2=3 \Rightarrow x=\pm\sqrt3$; \ $x^2=4 \Rightarrow x=\pm2$.
    \item Biquadratica, pongo $t=x^2$: \ $t^2-13t+36=0 \Rightarrow (t-4)(t-9)=0$. \\
    $x=\pm2$ e $x=\pm3$.
    \item \[\begin{aligned}
    x^2(x^2-3x-4) &= 0 && \textcolor{gray}{\text{raccolgo $x^2$}}\\
    x^2(x-4)(x+1) &= 0
    \end{aligned}\]
    $x=-1,\ 0,\ 4$.
    \item \[\begin{aligned}
    2x(x^2-4x+3) &= 0 && \textcolor{gray}{\text{raccolgo $2x$}}\\
    2x(x-1)(x-3) &= 0
    \end{aligned}\]
    $x=0,\ 1,\ 3$.
    \item \[\begin{aligned}
    x^2(x^2-x-6) &= 0 && \textcolor{gray}{\text{raccolgo $x^2$}}\\
    x^2(x-3)(x+2) &= 0
    \end{aligned}\]
    $x=-2,\ 0,\ 3$.
    \item Cerco una radice tra i divisori di $8$: \ $P(1)=1-7+14-8=0$, quindi $(x-1)$ è un fattore. \\
    Con la regola di Ruffini divido per $(x-1)$:
    \[\begin{aligned}
    (x-1)(x^2-6x+8) &= 0\\
    (x-1)(x-2)(x-4) &= 0 && \textcolor{gray}{\text{somma $-6$, prodotto $8$}}
    \end{aligned}\]
    $x=1,\ 2,\ 4$.
\end{enumerate}
\end{explanation}
\end{question}

%----------------------------------------------------%

\begin{question}[ID=4.5, subtitle={Equazioni di grado superiore al secondo con scomposizioni complesse}, difficulty=difficile, tags={equazioni di grado superiore}]
Risolvi le seguenti equazioni:
\begin{multicols}{2}
\begin{enumerate}[leftmargin=3cm]
    \item $x^4 - 2x^3 - 7x^2 + 8x + 12 = 0$
    \item $x^4 + x^3 - 7x^2 - x + 6 = 0$
    \item $x^3 - 6x^2 + 11x - 6 = 0$
    \item $x^4 - 5x^3 + 5x^2 + 5x - 6 = 0$
    \item $x^4 - 10x^2 + 9 = 0$
    \item $x^6 - 13x^3 + 36 = 0$
\end{enumerate}
\end{multicols}
\begin{solution}
\begin{enumerate}
    \item $x=-2,\ -1,\ 2,\ 3$
    \item $x=-3,\ -1,\ 1,\ 2$
    \item $x=1,\ 2,\ 3$
    \item $x=-1,\ 1,\ 2,\ 3$
    \item $x=-3,\ -1,\ 1,\ 3$
    \item $x=\sqrt[3]{4},\ \sqrt[3]{9}$
\end{enumerate}
\end{solution}
\begin{explanation}
Per i polinomi di grado alto si cercano le radici intere tra i divisori del termine noto (sostituendo nel polinomio) e si scompone con la regola di Ruffini.
\begin{enumerate}
    \item Divisori di $12$: provo $x=-1$: $P(-1)=1+2-7-8+12=0$; provo $x=2$: $P(2)=16-16-28+16+12=0$. \\
    Dividendo due volte con Ruffini:
    \[\begin{aligned}
    (x+1)(x-2)(x^2-x-6) &= 0\\
    (x+1)(x-2)(x-3)(x+2) &= 0 && \textcolor{gray}{\text{somma $-1$, prodotto $-6$}}
    \end{aligned}\]
    $x=-2,\ -1,\ 2,\ 3$.
    \item $P(1)=1+1-7-1+6=0$ e $P(-1)=1-1-7+1+6=0$. Ruffini due volte:
    \[\begin{aligned}
    (x-1)(x+1)(x^2+x-6) &= 0\\
    (x-1)(x+1)(x+3)(x-2) &= 0
    \end{aligned}\]
    $x=-3,\ -1,\ 1,\ 2$.
    \item $P(1)=1-6+11-6=0$. Ruffini:
    \[\begin{aligned}
    (x-1)(x^2-5x+6) &= 0\\
    (x-1)(x-2)(x-3) &= 0
    \end{aligned}\]
    $x=1,\ 2,\ 3$.
    \item $P(1)=1-5+5+5-6=0$ e $P(-1)=1+5+5-5-6=0$. Ruffini due volte:
    \[\begin{aligned}
    (x-1)(x+1)(x^2-5x+6) &= 0\\
    (x-1)(x+1)(x-2)(x-3) &= 0
    \end{aligned}\]
    $x=-1,\ 1,\ 2,\ 3$.
    \item Biquadratica, pongo $t=x^2$: \ $t^2-10t+9=0 \Rightarrow t=1,\ t=9$, quindi $x=\pm1,\ \pm3$.
    \item È un'equazione \textbf{trinomia}: pongo $t=x^3$.
    \[\begin{aligned}
    t^2-13t+36 &= 0\\
    (t-4)(t-9) &= 0 && \textcolor{gray}{\text{somma $13$, prodotto $36$}}
    \end{aligned}\]
    Torno alla $x$: $x^3=4 \Rightarrow x=\sqrt[3]{4}$; \ $x^3=9 \Rightarrow x=\sqrt[3]{9}$.
\end{enumerate}
\end{explanation}
\end{question}


%====================================================%
% SEZIONE 5: Legge dell'annullamento del prodotto
%====================================================%

\begin{question}[ID=5.1, subtitle={Legge dell'annullamento del prodotto}, difficulty=facile, tags={annullamento del prodotto}]
Risolvi le seguenti equazioni applicando la legge dell'annullamento del prodotto:
\begin{multicols}{2}
\begin{enumerate}[leftmargin=3cm]
    \item $(x-3)(x+2)=0$
    \item $(x+5)(x-1)=0$
    \item $x(x-7)=0$
    \item $(2x-4)(x+3)=0$
    \item $(3x+6)(x-2)=0$
    \item $(x-4)(x+4)=0$
\end{enumerate}
\end{multicols}
\begin{solution}
\begin{enumerate}
    \item $x=3,\ -2$
    \item $x=-5,\ 1$
    \item $x=0,\ 7$
    \item $x=2,\ -3$
    \item $x=-2,\ 2$
    \item $x=-4,\ 4$
\end{enumerate}
\end{solution}
\begin{explanation}
\textbf{Legge di annullamento del prodotto}: un prodotto vale zero se e solo se almeno uno dei suoi fattori vale zero. Si pone quindi uguale a zero ciascun fattore e si risolvono le equazioni ottenute.
\begin{enumerate}
    \item $x-3=0 \Rightarrow x=3$; \quad $x+2=0 \Rightarrow x=-2$.
    \item $x+5=0 \Rightarrow x=-5$; \quad $x-1=0 \Rightarrow x=1$.
    \item $x=0$; \quad $x-7=0 \Rightarrow x=7$.
    \item $2x-4=0 \Rightarrow 2x=4 \Rightarrow x=2$; \quad $x+3=0 \Rightarrow x=-3$.
    \item $3x+6=0 \Rightarrow 3x=-6 \Rightarrow x=-2$; \quad $x-2=0 \Rightarrow x=2$.
    \item $x-4=0 \Rightarrow x=4$; \quad $x+4=0 \Rightarrow x=-4$.
\end{enumerate}
\end{explanation}
\end{question}

%----------------------------------------------------%

\begin{question}[ID=5.2, subtitle={Legge dell'annullamento del prodotto}, difficulty=facile, tags={annullamento del prodotto}]
Risolvi le seguenti equazioni applicando la legge dell'annullamento del prodotto:
\begin{multicols}{2}
\begin{enumerate}[leftmargin=3cm]
    \item $(2x+1)(x-5)=0$
    \item $(3x-9)(x+4)=0$
    \item $(x-2)(4x+8)=0$
    \item $(5x-10)(2x+6)=0$
    \item $(x+1)(3x-12)=0$
    \item $(2x-6)(4x+1)=0$
\end{enumerate}
\end{multicols}
\begin{solution}
\begin{enumerate}
    \item $x=-\tfrac12,\ 5$
    \item $x=3,\ -4$
    \item $x=2,\ -2$
    \item $x=2,\ -3$
    \item $x=-1,\ 4$
    \item $x=3,\ -\tfrac14$
\end{enumerate}
\end{solution}
\begin{explanation}
Si pone uguale a zero ciascun fattore.
\begin{enumerate}
    \item $2x+1=0 \Rightarrow x=-\tfrac12$; \quad $x-5=0 \Rightarrow x=5$.
    \item $3x-9=0 \Rightarrow x=3$; \quad $x+4=0 \Rightarrow x=-4$.
    \item $x-2=0 \Rightarrow x=2$; \quad $4x+8=0 \Rightarrow x=-2$.
    \item $5x-10=0 \Rightarrow x=2$; \quad $2x+6=0 \Rightarrow x=-3$.
    \item $x+1=0 \Rightarrow x=-1$; \quad $3x-12=0 \Rightarrow x=4$.
    \item $2x-6=0 \Rightarrow x=3$; \quad $4x+1=0 \Rightarrow x=-\tfrac14$.
\end{enumerate}
\end{explanation}
\end{question}

%----------------------------------------------------%

\begin{question}[ID=5.3, subtitle={Legge dell'annullamento del prodotto con fattori di secondo grado}, difficulty=media, tags={annullamento del prodotto}]
Risolvi le seguenti equazioni applicando la legge dell'annullamento del prodotto:
\begin{multicols}{2}
\begin{enumerate}[leftmargin=3cm]
    \item $x(x-2)(x+3)=0$
    \item $(x-1)(x+2)(x-4)=0$
    \item $(2x-3)(x^2-4)=0$
    \item $(x+1)(x^2-9)=0$
    \item $(3x+2)(x^2-5x+6)=0$
    \item $(x-2)(x^2+x-6)=0$
\end{enumerate}
\end{multicols}
\begin{solution}
\begin{enumerate}
    \item $x=-3,\ 0,\ 2$
    \item $x=-2,\ 1,\ 4$
    \item $x=-2,\ \tfrac32,\ 2$
    \item $x=-3,\ -1,\ 3$
    \item $x=-\tfrac23,\ 2,\ 3$
    \item $x=-3,\ 2$
\end{enumerate}
\end{solution}
\begin{explanation}
Si pone uguale a zero ogni fattore; i fattori di secondo grado si risolvono (o si scompongono) a parte.
\begin{enumerate}
    \item $x=0$; \quad $x-2=0 \Rightarrow x=2$; \quad $x+3=0 \Rightarrow x=-3$.
    \item $x=1$, \ $x=-2$, \ $x=4$.
    \item $2x-3=0 \Rightarrow x=\tfrac32$; \quad $x^2-4=0 \Rightarrow x^2=4 \Rightarrow x=\pm2$.
    \item $x+1=0 \Rightarrow x=-1$; \quad $x^2-9=0 \Rightarrow x=\pm3$.
    \item $3x+2=0 \Rightarrow x=-\tfrac23$; \\
    $x^2-5x+6=0$: due numeri con somma $-5$ e prodotto $6$ sono $-2$ e $-3$, quindi $(x-2)(x-3)=0 \Rightarrow x=2,\ 3$.
    \item $x-2=0 \Rightarrow x=2$; \\
    $x^2+x-6=(x+3)(x-2)=0 \Rightarrow x=-3$ oppure $x=2$. \\
    $x=2$ compare due volte (soluzione doppia): le soluzioni distinte sono $x=-3$ e $x=2$.
\end{enumerate}
\end{explanation}
\end{question}

%----------------------------------------------------%

\begin{question}[ID=5.4, subtitle={Legge dell'annullamento del prodotto e scomposizione}, difficulty=media, tags={annullamento del prodotto}]
Risolvi le seguenti equazioni applicando la legge dell'annullamento del prodotto:
\begin{multicols}{2}
\begin{enumerate}[leftmargin=3cm]
    \item $x^2 - 5x + 6 = 0$
    \item $x^2 - 9 = 0$
    \item $x^3 - 4x^2 + 3x = 0$
    \item $x^3 + 2x^2 - 8x = 0$
    \item $x^3 - x^2 - 6x = 0$
    \item $x^4 - 5x^2 + 4 = 0$
\end{enumerate}
\end{multicols}
\begin{solution}
\begin{enumerate}
    \item $x=2,\ 3$
    \item $x=-3,\ 3$
    \item $x=0,\ 1,\ 3$
    \item $x=-4,\ 0,\ 2$
    \item $x=-2,\ 0,\ 3$
    \item $x=-2,\ -1,\ 1,\ 2$
\end{enumerate}
\end{solution}
\begin{explanation}
Prima si scompone il polinomio in fattori, poi si applica la legge di annullamento del prodotto.
\begin{enumerate}
    \item Trinomio: somma $-5$, prodotto $6$: \ $(x-2)(x-3)=0 \Rightarrow x=2,\ 3$.
    \item Differenza di quadrati: \ $(x-3)(x+3)=0 \Rightarrow x=\pm3$.
    \item \[\begin{aligned}
    x(x^2-4x+3) &= 0 && \textcolor{gray}{\text{raccolgo $x$}}\\
    x(x-1)(x-3) &= 0 && \textcolor{gray}{\text{somma $-4$, prodotto $3$}}
    \end{aligned}\]
    $x=0,\ 1,\ 3$.
    \item \[\begin{aligned}
    x(x^2+2x-8) &= 0 && \textcolor{gray}{\text{raccolgo $x$}}\\
    x(x+4)(x-2) &= 0 && \textcolor{gray}{\text{somma $2$, prodotto $-8$}}
    \end{aligned}\]
    $x=-4,\ 0,\ 2$.
    \item \[\begin{aligned}
    x(x^2-x-6) &= 0 && \textcolor{gray}{\text{raccolgo $x$}}\\
    x(x-3)(x+2) &= 0 && \textcolor{gray}{\text{somma $-1$, prodotto $-6$}}
    \end{aligned}\]
    $x=-2,\ 0,\ 3$.
    \item \[\begin{aligned}
    (x^2-1)(x^2-4) &= 0 && \textcolor{gray}{\text{trinomio in $x^2$: somma $-5$, prodotto $4$}}\\
    (x-1)(x+1)(x-2)(x+2) &= 0 && \textcolor{gray}{\text{differenze di quadrati}}
    \end{aligned}\]
    $x=-2,\ -1,\ 1,\ 2$.
\end{enumerate}
\end{explanation}
\end{question}

%----------------------------------------------------%

\begin{question}[ID=5.5, subtitle={Legge dell'annullamento del prodotto con più fattori}, difficulty=difficile, tags={annullamento del prodotto}]
Risolvi le seguenti equazioni applicando la legge dell'annullamento del prodotto e le opportune scomposizioni:
\begin{multicols}{2}
\begin{enumerate}[leftmargin=3cm]
    \item $(x^2-4)(x^2-9)=0$
    \item $(x^2-5x+6)(x^2-1)=0$
    \item $x(x-1)(x^2-4)=0$
    \item $(x^2-2x)(x^2-9)=0$
    \item $(x^2-4x+3)(x^2-16)=0$
    \item $(x^2-3x+2)(x^2+x-6)=0$
\end{enumerate}
\end{multicols}
\begin{solution}
\begin{enumerate}
    \item $x=-3,\ -2,\ 2,\ 3$
    \item $x=-1,\ 1,\ 2,\ 3$
    \item $x=-2,\ 0,\ 1,\ 2$
    \item $x=-3,\ 0,\ 2,\ 3$
    \item $x=-4,\ 1,\ 3,\ 4$
    \item $x=-3,\ 1,\ 2$
\end{enumerate}
\end{solution}
\begin{explanation}
Si scompone ogni fattore di secondo grado e si pone uguale a zero ciascun fattore di primo grado.
\begin{enumerate}
    \item $(x-2)(x+2)(x-3)(x+3)=0$ \ (differenze di quadrati) $\Rightarrow x=-3,\ -2,\ 2,\ 3$.
    \item $x^2-5x+6=(x-2)(x-3)$ e $x^2-1=(x-1)(x+1)$: \\
    $(x-2)(x-3)(x-1)(x+1)=0 \Rightarrow x=-1,\ 1,\ 2,\ 3$.
    \item $x^2-4=(x-2)(x+2)$: \\
    $x(x-1)(x-2)(x+2)=0 \Rightarrow x=-2,\ 0,\ 1,\ 2$.
    \item $x^2-2x=x(x-2)$ e $x^2-9=(x-3)(x+3)$: \\
    $x(x-2)(x-3)(x+3)=0 \Rightarrow x=-3,\ 0,\ 2,\ 3$.
    \item $x^2-4x+3=(x-1)(x-3)$ e $x^2-16=(x-4)(x+4)$: \\
    $(x-1)(x-3)(x-4)(x+4)=0 \Rightarrow x=-4,\ 1,\ 3,\ 4$.
    \item $x^2-3x+2=(x-1)(x-2)$ e $x^2+x-6=(x+3)(x-2)$: \\
    $(x-1)(x-2)^2(x+3)=0 \Rightarrow x=-3,\ 1,\ 2$ \ ($x=2$ è soluzione doppia).
\end{enumerate}
\end{explanation}
\end{question}